\section{Finite numerical estimates}\label{app:numeric}

We specify the finite quantities that determine the two constants
in Lemma~\ref{lem:numerical-endpoints}. The construction separates
integral enclosures, the finite defect recurrence, and the final
rational subtraction, so that each contribution has an explicit
mathematical definition and direction.

\subsection{Supplementary data}
The finite data are supplied in \emph{Theorem (1+1.75) on the Goldbach
Conjecture --- Numerical Supplement}, version 20261106~\cite{ChanSupplement}.
They comprise $106$ affine transitions, $200$-cell defect arrays at
depths $248$, $312$, and $430$, and $17572$ paired rectangles.
Rational quantities in \filepath{finite-construction-data.json} are
represented as quotients of integers. The fields
\texttt{upper\_saving\_bounds}, \texttt{lower\_root\_bounds}, and
\texttt{lower\_integral\_bounds} give directed rational endpoints.
The map \filepath{certificate-map.json} identifies their source
certificates and the corresponding domains, subdivisions, and interval
witnesses. Table~\ref{num:certificate-table} gives the correspondence
between the mathematical estimates and their certificate fields.

\hypertarget{num:verification-scope}{}
\paragraph{Verification commands and inputs.}
Run the commands below from the supplement root with Python 3.11 or later,
with assertions enabled. The sixth-rank computations also use NumPy.
The rows distinguish reconstruction from saved certificates, exact checks of
the finite construction, and recalculation of specified witness bounds.
All output files requested with \texttt{-{}-output} or \texttt{-{}-output-dir}
are written outside the supplement.

\begingroup
\renewcommand{\LTcaptype}{}
\stepcounter{LT@tables}
\small
\setlength{\tabcolsep}{4pt}
\begin{longtable}{@{}>{\raggedright\arraybackslash}p{.40\textwidth}>{\raggedright\arraybackslash}p{.29\textwidth}>{\raggedright\arraybackslash}p{.25\textwidth}@{}}
\toprule Entry point & Quantities checked or recomputed & Numerical inputs retained \\
\midrule
\endhead
\bottomrule
\endfoot
\filepath{verify-package.py} & File set, sizes and recorded SHA-256 identities. & The versioned \filepath{MANIFEST.json}.\\[3pt]
\filepath{finite-data/reproduce-data.py} \texttt{-{}-check} & Exact extraction of finite data and field correspondence from source certificates. & The supplied source-certificate values.\\[3pt]
\filepath{finite-data/verify-finite-data.py} & Finite defect recurrence; paired-cell geometry, masses and loss proxy; $\Gamma_5$ decompositions; density sums and strict scalar margins. & Action coefficients and integral endpoints identified by \filepath{certificate-map.json}.\\[3pt]
\filepath{replay-records.py} & Ordered-cover keys, stored integer-floor sums and the cost identity. & Stored floors in the ordered witness.\\[3pt]
\filepath{replay-ordered-cover.py} & The selected-method floor on every row of the grid-48 base cover; use \texttt{-{}-grid 32} for its predecessor. & The recorded method choices and generation-248 defect table.\\[3pt]
\filepath{verify-sixth-domains.py} & Affine sieve-argument ranges on the saved polyhedral cells. & Recorded vertices, coordinate maps and caps.\\[3pt]
\filepath{replay-sixth-rank.py} & Grid-400 sixth-rank cube contributions and their sum, together with the preceding polyhedral range checks. & The supplied cube witness and polyhedral records.\\
\end{longtable}
\endgroup

The recurrence and paired-cell checks use the action and integral inputs
listed in Table~\ref{num:certificate-table}. Their primitive enclosure
obligations remain those of the corresponding integrals and policies in
this appendix and Appendices~\ref{app:upper}--\ref{app:lower}.
The \texttt{enclosure\_obligations} entries of \filepath{certificate-map.json}
list the active assertions, complete records and supplied algorithm sources
for each primitive family. Witness generation is described in the package
README; in particular \filepath{witnesses/generator-0041.py} generates the ordered
base cover and \filepath{witnesses/generator-0045.py} the sixth-rank cube witness.
The generation-248 profile and other numerical inputs retained by a replay
are checked through their own entries, so that each verification has an
explicit set of inputs. The package README gives the command options; the
fixed-point arithmetic for the finite data is specified below.

\subsection{Exact scalar endpoints}
The exact values used in the upper density estimate are
\begin{align}
 E_{\rm up}&=\frac{9979702576710911234850619807816847864907359959}
 {1475913550962201227946780619099353215139840000},
 \label{eq:current-E}\\
 u&=\frac{145007961227048312210593427530128000784681}
 {11807308407697609823574244952794825721118720000}.
 \label{eq:current-u}
\end{align}
For the lower construction, before the paired-fibre increment, put
\begin{align}
 R_{\rm base}&=
 \frac{1345575129530003739571681432301188857312754737}
 {685738127494661151198444592354557952000000000},
 \label{eq:base-R}\\
 B_{\rm base}&=
 \frac{2504748318091942808229694219379799307550987}
 {5485905019957289209587556738836463616000000000}.
 \label{eq:base-B}
\end{align}
The two copies of $u$ and of $B_{\rm base}$ pay for distinct finite
approximation and coefficient-realization errors. Quadrature errors
are already contained in the oriented integral endpoints. The ancillary
fields \texttt{scalar\_bounds} contain the resulting exact sums and
strict surpluses.

\subsection{Directed arithmetic}
All rounding used to prove an upper bound is outward, and all rounding
used to prove a lower bound is inward.  More explicitly, at denominator
$\mathsf Q=2^{b_{\rm fp}}$ a rational interval $[x,y]$ is represented by
$[\lfloor \mathsf Q x\rfloor,\lceil \mathsf Q y\rceil]$.  Addition and subtraction
are exact at the integer level.  For multiplication the four endpoint
products are enclosed by their minimum and maximum and rounded in the
indicated directions after division by $\mathsf Q$.  Division is performed
only after verifying that the denominator interval is positive.

For logarithms put $z=(x-1)/(x+1)$.  If $1\le x\le2$, then
\begin{equation}
 2\sum_{j=0}^{n-1}\frac{z^{2j+1}}{2j+1}
 \le\log x\le
 2\sum_{j=0}^{n-1}\frac{z^{2j+1}}{2j+1}
 +\frac{2z^{2n+1}}{(2n+1)(1-z^2)}.
 \label{num:log-enclosure}
\end{equation}
This follows by summing the positive power series for
$2\operatorname{arctanh}z$ and bounding the remaining denominators
below by $2n+1$.  For $n=96$ the remainder is at most
$9/(4\cdot193\cdot3^{193})$.  Range reduction and a separate
enclosure of $\log2$ extend the formula to every positive rational
argument used here.  The paired-cell calculation uses denominator
$2^{224}$ for the series and denominator $2^{160}$ for its returned
intervals.  Thus no floating-point evaluation of a logarithm determines
the direction of an inequality.

\subsection{Enclosing delayed functions and integrals}
It is convenient to propagate the numerators
\[
 A_+(s)=\frac{sF(s)}{2e^\gamma},\qquad
 A_-(s)=\frac{sf(s)}{2e^\gamma}.
\]
Their initial branches are $A_+(s)=1$ for $1\le s\le3$ and
$A_-(s)=0$ for $s\le2$, and subsequent branches satisfy
\[
 A_+'(s)=\frac{A_-(s-1)}{s-1},\qquad
 A_-'(s)=\frac{A_+(s-1)}{s-1}.
\]
The equation for a branch involves only the preceding unit interval.
One therefore obtains rational interval enclosures successively on a
rational mesh, splitting at every integer breakpoint.  Lower and upper
integral bounds follow from the corresponding interval bounds for the
integrand.  Chord bounds and curvature bounds sharpen these enclosures
without changing their direction.  An arbitrary query is enclosed on
its entire support, including both adjoining cells when an endpoint
is a mesh knot.  The analogous construction for $s\omega(s)$ uses
$(s\omega(s))'=\omega(s-1)$ and its initial branch
$\omega(s)=1/s$ on $[1,2]$.

The chord integration can be implemented without numerical
quadrature.  On $[j/n,(j+1)/n]$ the integral of the affine
endpoint interpolant of a numerator $V$, divided by its argument,
is
\[
 \bigl((j+1)\log(1+1/j)-1\bigr)V(j/n)
 +\bigl(1-j\log(1+1/j)\bigr)V((j+1)/n).
\]
Both weights are positive because they are integrals of
nonnegative affine basis functions against $du/u$.
The monotonicities of $F$ and $f$ imply that $A_+$ is convex
and $A_-$ is concave on their nonconstant branches.
Their numerator curvature on a delayed cell is bounded
in absolute value by $1/(j/n-1)$.  Thus the integrated
chord error is at most
\[
 \frac{1}{12nj(j-n)}\qquad(j>n).
\]
This quantity is subtracted from a chord lower estimate
or added to a chord upper estimate as appropriate.
The weights themselves have alternating rational series
\[
 (j+1)\log(1+1/j)-1
  =\sum_{k\ge1}\frac{(-1)^{k+1}}{k(k+1)j^k},\qquad
 1-j\log(1+1/j)
  =\sum_{k\ge1}\frac{(-1)^{k+1}}{(k+1)j^k}.
\]
Their fourth and fifth partial sums enclose the exact
weights.  This supplies rational enclosures even where
direct subtraction of nearly equal logarithms would be
numerically inconvenient.

The finite integrations use rational boxes or rational polytopes.
For a polytope given by affine inequalities, every possible active
basis is solved exactly, and each resulting vertex is tested against
every original inequality.  Boundedness and the face incidences are
verified before a pulling triangulation is formed.  Polynomial
moments are then exact: on a $d$-simplex $T$ with barycentric
coordinates $\lambda_0,\ldots,\lambda_d$,
\begin{equation}
 \int_T\prod_{j=0}^d\lambda_j^{\alpha_j}\,dx
 =d!\operatorname{vol}(T)
   \frac{\prod_{j=0}^d\alpha_j!}{(d+\sum_j\alpha_j)!}.
 \label{num:simplex-moment}
\end{equation}
Rational lower or upper bounds for the remaining kernel factors,
together with these moments and \eqref{num:log-enclosure}, give the
required integral enclosure.  A boundary box in a negative-cost upper
integral is retained with an upper bound for its whole contribution.
In a positive saving, a box may instead be omitted after establishing
that its omitted contribution is nonnegative.  These two rules have
different signs and are used separately.

\subsection{The upper density}
Write
\[
 U_0=\frac4\Theta+I_1+I_{3,c}+I_{3,t},\qquad
 E_{\rm up}=U_0-\sum_{j=1}^{22}G_j.
\]
Here $I_{3,c}$ is the integral over the region with last coordinate
at least $1/100$, and $I_{3,t}$ bounds its entire complementary
tail.  Put $D=(F-f)/(2e^\gamma)$.  Its delay equation gives
\[
 D'(s)=-\frac{D(s)+D(s-1)}s\le-\frac{2D(s)}s
 \quad(s\ge3),
\]
using positivity and monotonicity of $D$.  Therefore
$D(s)\le D(10)(10/s)^2$ for $s\ge10$.  If
$\delta=1/100$ and $\rho_0=19/96$, the upper-tree cap
inequalities give a residual exponent at least $\rho_0$.
For $0<z<\delta$ its delayed argument is at least
$(\rho_0-z)/z\ge10$.  Enlarging the two remaining ordered
coordinates to $z\le y\le x\le\rho_0$ yields the complete
tail bound
\begin{equation}
 I_{3,t}\le
 \frac{200D(10)\delta}{(\rho_0-\delta)^2}
 \left\{\log^2\frac{\rho_0}{\delta}
        +2\log\frac{\rho_0}{\delta}+2\right\}.
 \label{num:complete-tail}
\end{equation}
Here the ordered $x,y$ integral is
$\tfrac12\log^2(\rho_0/z)$, and the remaining integral
over $0<z<\delta$ is elementary.  In particular the
logarithmic singularity at $z=0$ is integrated rather
than discarded.  The first integral is
enclosed by a positive series with its geometric remainder.
The compact third integral is covered by a disjoint partition,
including conservative bounds on its boundary boxes.

Table~\ref{num:upper-table} gives readable lower bounds for the
nonnegative savings $G_j$: its integer entries satisfy
$G_j\ge m_j10^{-14}$. The full rational bounds are in the
supplementary data. Successive changes on one region are measured
from their respective parents. In particular, the last two rows
give only the curvature and mixed increments over the preceding
bound for that integral.

\begin{longtable}{@{}r p{.62\textwidth} r@{}}
\caption{Rational lower bounds $10^{-14}m_j$ for the upper-density savings.}
\label{num:upper-table}\\
\toprule
$j$ & Saving & $m_j$\\
\midrule
\endfirsthead
\toprule $j$ & Saving & $m_j$\\\midrule
\endhead
1 & Fourth-rank cap replacement & 3507783546602\\
2 & Fifth-rank parent subdivision & 1005163763573\\
3 & Disjoint fifth-rank region & 136378695993\\
4 & Sixth-rank cap replacement & 96281257455\\
5 & Second-rank fallback above $\mu_0$ & 4209783419\\
6 & Third-rank fallback above $\mu_0$ & 3908979560\\
7 & Sixth-rank logarithmic moments & 1525398138\\
8 & Sixth-rank descending region & 187855687\\
9 & Sixth-rank ascending cap difference & 387157212\\
10 & Sixth-rank descending cap difference & 236377744\\
11 & Sixth-rank clipped lower window & 846201484\\
12 & Sixth-rank descending lower window & 500857555\\
13 & Sixth-rank local cap corrections & 50721206\\
14 & Sixth-rank nonsingleton regions & 11743733\\
15 & Eighth-rank radial bound & 94105689\\
16 & Eighth-rank ascending regions & 1543640141\\
17 & Eighth-rank descending regions & 590010094\\
18 & Eighth-rank unequal-parent moments & 1304343681\\
19 & Eighth-rank third-coordinate moments & 357686854\\
20 & Eighth-rank prefix-dependent moments & 37235787\\
21 & Eighth-rank curvature increment & 97705\\
22 & Eighth-rank mixed increment & 607514\\
\bottomrule
\end{longtable}

The upper enclosures use $3699127$ fourth-rank boxes,
$292633$ fifth-rank parents subdivided into $18728512$ children,
$1763420$ boxes in the disjoint fifth-rank region, and $632104$
sixth-rank cubes. The eighth-rank regions have disjoint families
of sizes $30$, $15$, $13$, and $29$. Across ranks the cap
differences telescope. On a parent region take the maximum of its
valid saving bounds, comparing its own bound with the complete
child sum.

An upper bound for the four base terms, in the same denominator
$10^{14}$, is respectively
\[
 668896321070235,\quad76020940667,\quad
 11960259849166,\quad12313775.
\]
Let $E_{\rm up,14}$ be their sum minus the sum of Table~\ref{num:upper-table},
divided by $10^{14}$.  Exact integer arithmetic gives
\[
 E_{\rm up}\le E_{\rm up,14},\qquad E_{\rm up,14}+2u<K,
\]
where $u$ is given in \eqref{eq:current-u}; indeed the last difference
is greater than $9.65\cdot10^{-8}$.  The unrounded rational sum
gives \eqref{eq:current-E}.  The two copies of $u$ pay separately
for finite approximation and coefficient realization.  Quadrature errors are included in the one-sided integral endpoints.

\subsection{The finite defect recurrence}
Put $\mathsf Q=2^{80}$ and
\[
 C_i=[1+i/50,1+(i+1)/50],\qquad 0\le i<200.
\]
All table entries in this subsection are integer numerators at
denominator $\mathsf Q$.  A four-step seed is obtained by starting at
$J_0=1/50$ and using
\[
 J_{n+1}=\frac{20249}{10^6}+\frac{13}{100}J_n.
\]
The initial defect and the implication represented by this recurrence
are supplied by Lemma~\ref{sv:finite-seed}, with its marked-state and
physical-cutoff conditions. The ancillary \texttt{seed\_certificate}
contains the separate ordinary and feedback lower endpoints used
there. Thus the arithmetic starts from an analytic saving valid on
the required family. Then $J_4=0.023273777353>2327/100000$. Initialize
\begin{equation}
 H_i^{(0)}=
 \begin{cases}
 \lfloor \mathsf Q\,2327/100000\rfloor,&1+(i+1)/50\le11/5,\\
 0,&\text{otherwise},
 \end{cases}
 \qquad h_i^{(0)}=0.
 \label{num:table-seed}
\end{equation}

For completeness the exact integer transition is specified here.
An action $\alpha$ consists of an integer $b_\alpha$, a divisor
$d_\alpha\in\{1,5\}$, point pairs $(j,m)$, nonnegative integer
weights $w_{\alpha j},v_{\alpha j}$, and a target set $T_\alpha$.
Its candidate value is
\begin{equation}
 b_\alpha+\left\lfloor\frac1{d_\alpha}
 \left\{\sum_{(j,m)}mH_j+
 \sum_j\left\lfloor\frac{w_{\alpha j}h_j}{\mathsf Q}\right\rfloor+
 \sum_j\left\lfloor\frac{v_{\alpha j}H_j}{\mathsf Q}\right\rfloor
 \right\}\right\rfloor.
 \label{num:action-map}
\end{equation}
The $70$ one-step actions have $d_\alpha=1$ and the $36$ two-step
actions have $d_\alpha=5$.  They are evaluated at the upper
constant $K_*=6.761767>K$.  The action parameters, weights,
constants and target sets are given individually in the
\texttt{actions} array.  The target condition for action parameter
$s_\alpha$ is $1+(i+1)/50\le s_\alpha$.

Let $\ell_j$ be the integer lower bound for
$\mathsf Q\log((1+(j+1)/50)/(1+j/50))$ in the data, and set
\[
 S_H(j)=\sum_{k=j}^{199}\left\lfloor\frac{H_k\ell_k}{\mathsf Q}\right\rfloor,
 \qquad
 S_h(j)=\sum_{k=j}^{199}\left\lfloor\frac{h_k\ell_k}{\mathsf Q}\right\rfloor.
\]
For $50\le i<200$ the Buchstab candidates are
$S_h(\max(50,i-49))$ for $H_i$ and $S_H(i-49)$ for $h_i$.
Each new entry is the maximum of its preceding value and all
applicable candidates.  Every right-hand side is evaluated using
only the preceding arrays.  This is a simultaneous finite
recurrence, with no assumption of convergence to a fixed point.

The analytic action inequalities and their closure under descendants
imply by induction that the rational arrays are lower bounds for
the defect functions on the whole closed cells.  Floors in
\eqref{num:action-map} preserve this direction because the
feedback weights are nonnegative.  Applying the map $430$ times
to \eqref{num:table-seed} gives the $400$ entries in
\texttt{defect\_tables.430}.  The arrays at depth $248$ are
bounded below by the corresponding finite iteration.  The
depth-$430$ arrays are used only in the two occurrences
$U_3,U_4$; the other designated occurrences use the depth-$248$
arrays.  Thus a change of table is charged only at its actual
occurrence in the lower functional.

\subsection{The integrals attached to the eleven counts}
\label{num:eleven-integrals}
The correspondence with the arithmetic terms $Y_j$ is fixed as follows.
Write $A_+(s)=sF(s)/(2e^\gamma)$,
$A_-(s)=sf(s)/(2e^\gamma)$, and put
$\nu=249/448$, $w_*=25/224$. The reference one-prefix terms are
\begin{align*}
 L_1&=\frac4{a_{\rm s}}\overline f(\nu/{a_{\rm s}}),\\
 L_2&=4\left\{(1/{h_{\rm s}}+1)\overline f\left(\frac{1+{h_{\rm s}}}{2{h_{\rm s}}}\right)
                 -\overline f(249/50)\right\}+H_{\rm rem},
\end{align*}
where $H_{\rm rem}$ is defined in~\eqref{num:removed-defect}.
No defect table is inserted at the unmarked root of $L_1$:
its ratio $1/\nu$ is less than $2$.

For the two-prefix lower terms put
\[
 \mathcal P_5=\{{a_{\rm s}}<u<t<{h_{\rm s}}\},\qquad
 \mathcal P_6=\{{a_{\rm s}}<u<{h_{\rm s}}<t<{d_{\rm s}}\}.
\]
On the enhanced part $t<\mu_0$, let
$V(t,u)$ be the maximum of the admitted candidates
$v(t),v(t+u),\Lambda_{\rm cap}-3{a_{\rm s}}/8$; on $t>\mu_0$ take $V=v(t)$.
Here admission retains the full and selected squares, the selected
largest slots and the physical prime ceiling, as in
Appendix~\ref{app:coefficients}. Set $L=V-t-u$. The lower integrals are
\begin{equation}
 L_j=4\iint_{\mathcal P_j}
        \frac{A_-(L/{a_{\rm s}})+h_{248,\rm adm}(t,u)}{tuL}\,du\,dt,
 \qquad j=5,6.
 \label{num:pair-lower-integrals}
\end{equation}
The defect term equals $h_{248}(L/{a_{\rm s}})$ only on a fully admitted state
with $2\le L/{a_{\rm s}}\le5$, and is zero otherwise. The ordinary lower term
is also zero when $L/{a_{\rm s}}<2$. Thus an uncertain defect lookup can be
discarded without changing an ordinary lower estimate.

For $m=248,312,430$, the one-prefix upper integrals are
\begin{equation}
 U_j^{(m)}=4\int_{I_j}
       \frac{A_+(s_j(t))-H_m(s_j(t))}{tL_j(t)}\,dt,
 \quad I_3=({a_{\rm s}},1/3),\quad I_4=({a_{\rm s}},{d_{\rm s}}),
 \label{num:one-prefix-upper-integrals}
\end{equation}
with $L_j,s_j$ given in~\eqref{num:table-change-integral} below.
The reference values in the eleven-term sum use $m=248$.
The two subsequent table improvements are inserted only through
\eqref{num:successive-table-gains}.

The last two costs use decreasing coordinates $x>y>z>w$ on
\[
 \mathcal D_{10}=\{{a_{\rm s}}<w<z<y<x<{h_{\rm s}}\},\qquad
 \mathcal D_{11}=\{{a_{\rm s}}<w<z<y<{h_{\rm s}}<x<{d_{\rm s}}\}.
\]
The integration measure is
\[
 {d_{\rm s}}\sigma_4=\frac{dx\,dy\,dz\,dw}{xyzw}.
\]
Put $S=x+y+z+w$. The switched option has kernel
$K_0\omega((1-S)/z)/z$. Each admitted ordinary branch with residual
exponent $L=\lambda-S$ supplies the kernel
$4\{1-H_{248}(\max(2,L/w))\}/L$; all its full and selected support conditions
must hold on the whole cell. The chosen cell majorant bounds the
minimum of these options from above. Its integral defines the upper
enclosure for $U_{10}+U_{11}$. The cell choices, residual intervals,
closed-cell defect minima and complete ordered cover are the data
of the four-prefix certificate. The switched option remains available
when an ordinary branch fails its support conditions. Since ${d_{\rm s}}+3{h_{\rm s}}<1$, these two
domains have no unit-quotient contribution.
The remaining switched costs $U_7,U_8,U_9$ are defined in
Section~\ref{num:cost-slope}.

\subsection{Definitions of the lower correction terms}
\label{num:lower-definitions}
We specify the correction terms before giving their numerical
combination. Put $\nu=249/448$, $w_*=25/224$ and write
$\overline F=F/(2e^\gamma)$, $\overline f=f/(2e^\gamma)$.
The difference between the ordinary start at ${h_{\rm s}}$ and the ordinary
part of the least-prime decomposition is
\begin{equation}
 B_{\rm ord}=\frac4{h_{\rm s}}\overline f(\nu/{h_{\rm s}})
 -4\left\{(1/{h_{\rm s}}+1)\overline f\left(\frac{1+{h_{\rm s}}}{2{h_{\rm s}}}\right)
                   -\overline f(249/50)\right\}.
 \label{num:ordinary-start-difference}
\end{equation}
The defect part removed from that decomposition is the whole integral
\begin{equation}
 H_{\rm rem}=4\int_{w_*}^{{h_{\rm s}}}
   \frac{H_{248}((1-t)/(2t))}{t(1-t)/2}\,dt.
 \label{num:removed-defect}
\end{equation}
The finite ordinary estimate at ${h_{\rm s}}$ is inserted only after this
entire defect term has been removed. On a defect cell $[s_l,s_r]$,
the contributing interval in~\eqref{num:removed-defect} is
$(w_*,{h_{\rm s}})\cap[1/(2s_r+1),1/(2s_l+1)]$; its primitive is
$2\log(t/(1-t))$, before multiplication by the factor $4$.

For $w={a_{\rm s}},{h_{\rm s}}$, let $C_3,C_4,C_5$ be the successive caps of the same
lower policy in Appendix~\ref{app:lower}. A cap is inherited or
selected from the finitely specified subproduct branches subject to
all full and selected-complement support conditions. In particular
$\nu\le C_3\le C_4\le C_5\le\Lambda_{\rm cap}$. For decreasing
$t_1,\ldots,t_r$, write $S_r=\sum_{j\le r}t_j$ and
\[
 {d_{\rm s}}\mu_r=\frac{dt_1\cdots dt_r}{t_1\cdots t_{r-1}t_r^2}.
\]
The three retained domains are
\begin{align*}
 \mathcal E_3(w)&=\{0<t_3<t_2<t_1<w:t_1+3t_2<\nu\},\\
 \mathcal E_4(w)&=\{(t_1,\ldots,t_4):(t_1,t_2,t_3)\in\mathcal E_3(w),
          \ 0<t_4<t_3,\ S_3+3t_4<C_3\},\\
 \mathcal E_5(w)&=\{(t_1,\ldots,t_5):(t_1,\ldots,t_4)\in\mathcal E_4(w),
                                      \ 0<t_5<t_4\}.
\end{align*}
Thus the three nonnegative density gains are
\begin{align}
 I_3(w)&=4\int_{\mathcal E_3(w)}
  \left\{\overline F\left(\frac{\nu-S_3}{t_3}\right)
        -\overline F\left(\frac{C_3-S_3}{t_3}\right)\right\}\,{d_{\rm s}}\mu_3,
                                      \label{num:third-lower-gain}\\
 I_4(w)&=4\int_{\mathcal E_4(w)}
  \left\{\overline f\left(\frac{C_4-S_4}{t_4}\right)
        -\overline f\left(\frac{C_3-S_4}{t_4}\right)\right\}\,{d_{\rm s}}\mu_4,
                                      \label{num:fourth-lower-gain}\\
 I_5(w)&=4\int_{\mathcal E_5(w)}
  \left\{\overline F\left(\frac{C_4-S_5}{t_5}\right)
        -\overline F\left(\frac{C_5-S_5}{t_5}\right)\right\}\,{d_{\rm s}}\mu_5.
                                      \label{num:fifth-lower-gain}
\end{align}
The fifth square follows from
$S_4+2t_5<S_3+3t_4<C_3\le C_4$.
The cap choices, branch labels and parent cubes are fixed jointly;
an integral evaluated with a different fourth cap cannot be appended
to~\eqref{num:fifth-lower-gain}. The rational certificates may integrate
only a subset of these nonnegative gains, so their endpoints are
lower bounds, rather than assertions that the selected boxes exhaust
the ideal domains.

For the fourth-node restoration, let $\mathcal R_{h_{\rm s}}$ be the disjoint
union of the rational slanted slabs specified by
\eqref{sv:restore-gates}. On each slab a cap $\lambda_4$ is chosen
before the fourth cube, and
$C_3\le S_3+3t_4<\lambda_4$, with the full and selected support
conditions retained. Its additional gain is
\begin{equation}
 T_{h_{\rm s}}=4\int_{\mathcal R_{h_{\rm s}}}
       \overline f\left(\frac{\lambda_4-S_4}{t_4}\right)\,{d_{\rm s}}\mu_4.
 \label{num:restoration-integral}
\end{equation}
The exact slab volume is~\eqref{sv:slab-volume}. This term enters the
Chen functional with coefficient $1/4$, and its domain is disjoint
from the other gains already counted.

\subsection{The fixed upper-cost slope and table changes}
\label{num:cost-slope}
The coefficient $A_{\rm cost}$ used below is a safe slope of a specified affine
upper enclosure of five costs. It need not be the derivative of a
minimum of kernels. To state its components, put $K_0=6.7619$ and
$A_+(s)=s\overline F(s)$. The three one-dimensional switched
costs are
\begin{align*}
 \frac{U_7}{K_0}&=\int_{d_{\rm s}}^{1/3}
  \frac{\log((1-2t)/t)+A_+(1/t-1)-1}{t(1-t)}\,dt,\\
 \frac{U_8}{K_0}&=\int_{a_{\rm s}}^{1/3}\frac{\log(2-3t)}{t(1-t)}\,dt,\\
 \frac{U_9}{K_0}&=\int_{h_{\rm s}}^{d_{\rm s}}\frac{\log((1-t-{d_{\rm s}})/{d_{\rm s}})}{t(1-t)}\,dt.
\end{align*}
The extra ordinary-function contribution in the first integral can
be evaluated without interpolation:
\[
 \int_{d_{\rm s}}^{1/4}\frac{A_+(1/t-1)-1}{t(1-t)}\,dt
 =\int_2^{1/{d_{\rm s}}-2}\frac{\log(u-1)}u
                  \log\frac{1-{d_{\rm s}}}{{d_{\rm s}}(u+1)}\,du.
\]
Let $\underline U_j$ denote the rational lower endpoints for these
$K_0$-linear costs from the fourth-derivative Simpson enclosures
specified in \texttt{five\_cost\_slope\_certificate}. The root cost
upper endpoints may use the separate midpoint enclosures; an upper
endpoint and a lower endpoint have different roles here.
Their contribution to a safe slope is
$A_{{\rm cost},789}=(2\underline U_7+\underline U_8+\underline U_9)/(4K_0)$.
For a retained four-prefix cell $\mathscr B$ outside the $40000$ priority
parents, write $A_{\mathscr B}$ for its assigned upper mass at $K_0$ and $s_{\mathscr B}$
for the supplied nonnegative affine slope. The required cellwise
conditions are
\begin{equation}
 0\le s_{\mathscr B}\le A_{\mathscr B}/K_0,\qquad A_{\mathscr B}\ge K_0\mathcal T_{\mathscr B},
 \label{num:affine-majorant-slope}
\end{equation}
where $\mathcal T_{\mathscr B}$ is the exact nonnegative switched mass. They give
\[
 A_{\mathscr B}-(K_0-K)s_{\mathscr B}\ge(K/K_0)A_{\mathscr B}\ge K\mathcal T_{\mathscr B}.
\]
No inequality $s_{\mathscr B}\le\mathcal T_{\mathscr B}$ is required.
At denominator $\mathsf Q=2^{80}$ the retained row slopes are
$s_{\mathscr B}=10000d_{\mathscr B}/\mathsf Q$. Their sums for $U_{10},U_{11}$ are
$1277985388931956822$ and $7976183963523684871$, respectively.
Consequently
\begin{equation}
 A_{\rm cost}=A_{{\rm cost},789}+
 \frac{10000(1277985388931956822+7976183963523684871)}{4\mathsf Q}.
 \label{num:five-cost-slope}
\end{equation}
The complete retained-cell inequalities are indexed by the
\texttt{feedback\_original\_slope} entry of the certificate map.

For two valid defect envelopes $H_m\le H_n$, define
\begin{equation}
 \mathcal T(H_m,H_n)=
 \int_{a_{\rm s}}^{1/3}\frac{(H_n-H_m)(s_3(t))}{tL_3(t)}\,dt
 +\int_{a_{\rm s}}^{d_{\rm s}}\frac{(H_n-H_m)(s_4(t))}{tL_4(t)}\,dt,
 \label{num:table-change-integral}
\end{equation}
where $s_j(t)=\max(2,L_j(t)/{a_{\rm s}})$, with zero contribution outside
$[1,5]$, and, for $e=57/224$,
\[
 L_3(t)=\begin{cases}
 (1-t)/2,&{a_{\rm s}}<t<\mu_0,\\
 57/64-5t/2,&\mu_0<t<e,\\
 1/2-t,&e<t<1/3.
 \end{cases}
\]
The function $L_4$ uses the first two branches restricted to $({a_{\rm s}},{d_{\rm s}})$.
The factors $4$ in the two upper counts cancel the single Chen
factor $1/4$. Thus
\begin{equation}
 T_{312}=\mathcal T(H_{248},H_{312}),\qquad
 T_{430}=\mathcal T(H_{312},H_{430}).
 \label{num:successive-table-gains}
\end{equation}
At a knot the lookup uses the smaller adjacent value; on the
positive-measure branch $s_j=2$ this means the minimum of cells $49$
and $50$. On an affine piece $L=A_0-B_0t$ the exact inverse-cell
primitive is
\[
 \int_l^r\frac{dt}{t(A_0-B_0t)}
 =\frac1{A_0}\log\frac{r(A_0-B_0l)}{l(A_0-B_0r)}.
\]
This determines each table-change interval and its rational logarithm
enclosure without numerical interpolation across a knot.

\subsection{Mathematical objects and certificate fields}
\label{num:certificate-chain}
Table~\ref{num:certificate-table} links the count estimates, their
integral bounds and the ancillary certificates. Paths in the last
column are targets in \filepath{certificate-map.json}; each target has
the source filename, exact JSON pointer, direction and any rational
extraction formula. The map also includes
the enclosure and coverage assertions for each primitive family and
the corresponding full-witness filenames.
The map's \texttt{domain\_constructions} entries identify the actual
$C_3,C_4,C_5$ region generators and their selected records, including
the two third-rank starts, the fourth-rank chord and oblique regions,
the fifth-rank child increments, and the slanted restoration.
For $U_{10},U_{11}$ they identify the generator
\filepath{witnesses/generator-0041.py} and the complete ordered witness
\filepath{witnesses/record-0638.json}.
The witness records specify the parent policy, method choices,
generation-$248$ profile, and full facet descriptions of these domains.

\begingroup
\small
\setlength{\tabcolsep}{4pt}
\begin{longtable}{@{}>{\raggedright\arraybackslash}p{.15\textwidth}>{\raggedright\arraybackslash}p{.43\textwidth}>{\raggedright\arraybackslash}p{.35\textwidth}@{}}
\caption{Definitions, directions and certificate targets.}
\label{num:certificate-table}\\
\toprule Object & Mathematical estimate & Certificate target\\
\midrule
\endfirsthead
\toprule Object & Mathematical estimate & Certificate target\\
\midrule
\endhead
$G_j$, $1\le j\le22$ &
The specified child--parent density differences in the upper tree;
nonnegative savings may use a subdomain. &
\filepath{/upper_saving_bounds/}\,$(j-1)$; lower endpoint.\\[3pt]
Finite actions &
\eqref{sv:one-step-operator}, \eqref{sv:two-step-base},
\eqref{sv:two-step-feedback}; complete $\Gamma_5$ with separate payment
on $D_5^-$. &
\filepath{/actions/}\,$k$; lower base and feedback weights;
\filepath{/gamma5_cost_decompositions}.\\[3pt]
Seed and tables &
Lemma~\ref{sv:finite-seed} and the simultaneous map
\eqref{num:action-map}; closed-cell estimates. &
\filepath{/seed_certificate}, \filepath{/defect_tables/}\,$m$;
lower endpoints, $m=248,312,430$.\\[3pt]
$L_i,U_i$ &
The integrals in Section~\ref{num:eleven-integrals}; complete retained
domain for each negative cost, lower estimate for each positive term. &
\filepath{/lower_root_bounds}; lower for $L_i$, upper for $U_i$.\\[3pt]
$B_{\rm ord},H_{\rm rem}$ &
\eqref{num:ordinary-start-difference}, \eqref{num:removed-defect};
remove the entire generation-$248$ defect. &
\filepath{/lower_auxiliary_bounds}; lower for $B_{\rm ord}$, upper for
$H_{\rm rem}$.\\[3pt]
$I_r(w),T_{h_{\rm s}}$ &
\eqref{num:third-lower-gain}--\eqref{num:restoration-integral}; one
common lower policy and disjoint restoration slabs. &
\filepath{/lower_integral_bounds} and
\filepath{/lower_auxiliary_bounds/restoration_L2_lower}; lower.\\[3pt]
$A_{\rm cost}$, $T_{312},T_{430}$ &
\eqref{num:affine-majorant-slope}--\eqref{num:successive-table-gains};
upper-cost slope and successive table changes. &
\filepath{/five_cost_slope_certificate},
\filepath{/lower_auxiliary_bounds}; safe lower coefficients.\\[3pt]
Paired deletion &
Assigned positive majorant and discarded lower estimator on
$\mathscr E_0\cup \mathscr E_1$; \eqref{num:paired-extra}, \eqref{num:saved-loss}. &
\filepath{/paired_rectangles}; exact geometry, mass intervals and lower
estimator integers.\\[3pt]
$E_{\rm up},R_{\rm low}$ and reserves &
The unique combinations in~\eqref{num:lower-sum} and the upper
density expression. &
\filepath{/scalar_bounds}; exact sums and strict surpluses.\\
\bottomrule
\end{longtable}
\endgroup

For each negative cost, the mapped witnesses specify a complete
cover of its retained domain. For a positive saving, they specify
the selected nonnegative subfamily on which its lower bound is
evaluated.

\subsection{The lower functional}
For the signed eleven-term expression, use a lower endpoint for
each positive term and an upper endpoint for each negative term.
The following integers, divided by $10^{14}$, give such endpoints:
\begin{center}
\begin{tabular}{@{}l r r@{}}
\toprule Term & Coefficient before division by $4$ & Endpoint numerator\\
\midrule
$L_1$ & $3$ & $1346406373358053$\\
$L_2$ & $1$ & $710262601191162$\\
$L_5$ & $1$ & $237450056438739$\\
$L_6$ & $1$ & $301984961684469$\\
$U_3$ & $-1$ & $2031242249325404$\\
$U_4$ & $-1$ & $1405868301028074$\\
$U_7$ & $-2$ & $138343028785273$\\
$U_8$ & $-1$ & $405707594761181$\\
$U_9$ & $-1$ & $325703365221150$\\
$U_{10}+U_{11}$ & $-1$ & $67568703438104$\\
\bottomrule
\end{tabular}
\end{center}
The exact rational endpoints, rather than these shortened ones,
define $C_{11}$.  The four-prefix minimum covers all
\[
 \binom{51}{4}+48\binom{50}{3}=1190700
\]
ordered rows, including the tied-coordinate simplices.
The priority subdivision replaces $40000$ parents by maxima
with their complete sums of $516779$ children.  Its extra
coefficient saving is
\[
 P=\frac{106520738437439293219}{604462909807314587353088}.
\]

The terms $B_{\rm ord},H_{\rm rem}$, $I_r(a_{\rm s}),I_r(h_{\rm s})$, $T_{h_{\rm s}}$ are
defined in~\eqref{num:ordinary-start-difference}--\eqref{num:restoration-integral}.
Their exact bounds are
specified by \texttt{lower\_auxiliary\_bounds} and
\texttt{lower\_integral\_bounds}.  At $K_0=6.7619$ put
\[
 C_{K_0}=C_{11}+P+\frac14\left(
 B_{\rm ord}-H_{\rm rem}
 +\sum_{r=3}^5\{3I_r({a_{\rm s}})+I_r({h_{\rm s}})\}+T_{h_{\rm s}}\right).
\]
The third-rank integrals cover both cutoff regions.  The
fourth-rank integrals retain all specified masks and the
disjoint additional bands.  The fifth-rank integration uses
$158251$ floors and $23662$ additional slab or tied cells on
the ${h_{\rm s}}$ side and $679$ integrals on the ${a_{\rm s}}$ side.  The slanted
restoration uses $98668$ boxes and $789344$ layers.  Only
nonnegative gains may omit boundary strips.  An asymmetric
kernel is bounded on its ordered simplex; no factorial
symmetry is assumed for such a kernel.

With $A_{\rm cost}$ as in~\eqref{num:five-cost-slope} and
$T_{312},T_{430}$ as in~\eqref{num:successive-table-gains}, the
initial lower endpoint is bounded below by
\begin{equation}
 C_{K_0}+A_{\rm cost}(K_0-K)+T_{312}+T_{430}
 +0.002095161+0.003799514-0.002372094.
 \label{num:lower-sum}
\end{equation}
Direct rational substitution is at least $R_{\rm base}$ in
\eqref{eq:base-R}.  The paired-cell refinement below adds
exactly $\Delta$ from \eqref{eq:current-refinement-delta}.
After allocating the two reserves in \eqref{eq:current-B},
the resulting lower endpoint exceeds $c$ by
\[
 \frac{13457260926863633771775678584711245753}
 {2742952509978644604793778369418231808000000000}>0.
\]
Together with the upper calculation, this proves the scalar
inequalities in Lemma~\ref{lem:numerical-endpoints} from the
stated finite enclosures.

\section{A paired-cell refinement}\label{app:refinement}

The paired rectangles are specified without a geometric
rounding convention as follows.  Put
\[
 t_i=\mu_0+i\frac{{d_{\rm s}}-\mu_0}{512},\qquad
 u_j={a_{\rm s}}+j\frac{{h_{\rm s}}-{a_{\rm s}}}{512}\qquad(0\le i,j\le512).
\]
Each row of \texttt{paired\_rectangles} specifies a family
$\mathscr E_0$ or $\mathscr E_1$ and the rectangle
$[t_i,t_{i+1}]\times[u_j,u_{j+1}]$.  The families have
$6091$ and $11481$ cells, respectively, with no common
cell.  Their boundaries have zero integration measure.
For a rectangle $\mathscr B=[t_l,t_r]\times[u_l,u_r]$, define
\[
 Z_{\mathscr B}=\min\left(u_l,\frac{3(1-t_r-u_r)}{13}\right).
\]
The rational data satisfy $Z_{\mathscr B}>{a_{\rm s}}$ and
$1-t_r-u_r-2Z_{\mathscr B}\ge(7/3)Z_{\mathscr B}$ in every row.  Hence
\[
 \frac{1-t-u-z-w}{z}\ge\frac73
 \qquad({a_{\rm s}}<w<z<Z_{\mathscr B},\ (t,u)\in \mathscr B).
\]

On each rectangle, Lemma~\ref{lem:current-omega-refinement} gives
$\omega\ge1/2$ throughout the fibre and $\omega\ge11/20$
on its specified inner part. The second estimate improves a
subset of the first contribution and is used as an increment.

The positive parent majorant can therefore be bounded below
by the contribution from $1/2$ on its whole fibre plus
$1/20$ on this inner fibre.  These are simultaneous
nonnegative contributions.  After integration and the
single Chen factor $1/4$, the additional contribution is
\begin{equation}
 G_{\mathscr B}=\frac K{80}\log\frac{t_r}{t_l}\log\frac{u_r}{u_l}
 \left\{\frac1{a_{\rm s}}-\frac{1+\log(Z_{\mathscr B}/{a_{\rm s}})}{Z_{\mathscr B}}\right\}.
 \label{num:paired-extra}
\end{equation}
Indeed the $w$ integral is $\log(z/{a_{\rm s}})$ and
\[
 \int_{a_{\rm s}}^{Z_{\mathscr B}}\frac{\log(z/{a_{\rm s}})}{z^2}\,dz
 =\frac1{a_{\rm s}}-\frac{1+\log(Z_{\mathscr B}/{a_{\rm s}})}{Z_{\mathscr B}}.
\]
The normalization agrees with~\eqref{eq:current-extra-fibre}.

Every row of the data contains rational lower and upper
endpoints for \eqref{num:paired-extra}, represented by
integer numerators at denominator $2^{160}$.  Applying
\eqref{num:log-enclosure} to its three logarithms and
summing the lower endpoints gives
\begin{equation}
 \sum_{\mathscr B\in \mathscr E_0\cup \mathscr E_1}G_{\mathscr B}>\frac{558583}{10^9}.
 \label{num:gross-bound}
\end{equation}
The sums on the two families are approximately
$1.9318091250\cdot10^{-4}$ and
$3.6540308174\cdot10^{-4}$, respectively. Only the rational
lower bound \eqref{num:gross-bound} is used.

\section{Discarded lower estimators and the final reserves}
\label{app:refinement-review}

The paired deletion removes a positive term together with its
assigned lower estimator. To bound the resulting loss, we give
both the exact finite sum of that estimator and a direct rational
upper enclosure.

On a cell write
\[
 R(t,u)=\frac{57}{64}-\frac52t-u,\qquad
 r=R(t_r,u_r),\quad R_0=R(t_l,u_l),\quad
 M=\log(t_r/t_l)\log(u_r/u_l).
\]
Cells of $\mathscr E_0$ have $R_0/{a_{\rm s}}\le2$ and zero assigned lower estimator.
There are $227$ cells of $\mathscr E_1$ with $r/{a_{\rm s}}<2$; their
assigned lower estimator is also zero.  On the remaining $11254$ cells,
$2\le r/{a_{\rm s}}\le R_0/{a_{\rm s}}\le9/4$.  Put
\[
 f_l=\frac{\log(r/{a_{\rm s}}-1)}{r/{a_{\rm s}}},\qquad
 f_h=\frac{\log(R_0/{a_{\rm s}}-1)}{R_0/{a_{\rm s}}},
\]
and
\[
 M_R=\frac{57}{64}M
 -\frac52(t_r-t_l)\log(u_r/u_l)
 -(u_r-u_l)\log(t_r/t_l).
\]
Since $r\le R(t,u)\le R_0$, the two assigned ordinary
estimators, in Chen-quarter units, are bounded above by
\begin{equation}
 \frac{f_lM}{{a_{\rm s}}},\qquad
 \frac{f_l(R_0M-M_R)+f_h(M_R-rM)}{a_{\rm s}(R_0-r)}.
 \label{num:ordinary-proxy}
\end{equation}
The coefficients of $f_l,f_h$ are nonnegative.
Replacing lower endpoints for these functions by their
exact values therefore increases the estimator.  The
maximum of the two expressions in
\eqref{num:ordinary-proxy} remains an upper bound for
the maximum of the two assigned alternatives.

Let $h_{\mathscr B}^{\min}$ be the minimum of the depth-$248$ lower table
entries over all closed cells meeting $[r/{a_{\rm s}},R_0/{a_{\rm s}}]$.
Both adjoining cells are included at a knot.  The exact
support is contained in the outward-rounded lookup
support of the assigned estimator, so $h_{\mathscr B}^{\min}$ is at least
the minimum used there.  Thus the additional term
$h_{\mathscr B}^{\min}M/R_0$ bounds its assigned defect estimator from above.
This comparison uses minima of finite lower estimators,
not an upper estimate for the mathematical function $h$.
Adding it to the maximum in
\eqref{num:ordinary-proxy} gives the required cell proxy.

The data record the assigned lower estimator on every cell as
an integer at denominator $\mathsf Q=2^{80}$.  Their sum is
$11290806845175652020270$.  After division by $4\mathsf Q$,
\begin{equation}
 L_{\mathscr E}
 =\frac{5645403422587826010135}
 {2417851639229258349412352}
 <\lambda_{\rm loss}:=\frac{466977}{200000000}.
 \label{num:saved-loss}
\end{equation}
An outward evaluation of the proxies above gives an
upper endpoint below the same $\lambda_{\rm loss}$.  For this
evaluation one may reduce logarithm arguments to
$[1,3/2]$, use the $320$-term alternating series for
$\log(1+x)$ and bound its remainder by the next term.
All its divisors have positive interval lower endpoints.
The resulting proxy sum is enclosed near
$0.0023348841363598131014744273487904$.

The retained whole $L_6$ endpoint before the deletion is
$3650774173156900914101685/\mathsf Q$.  Its nonenhanced half
has the cell decomposition just used.  Consequently
subtracting $4\lambda_{\rm loss}$ from that endpoint bounds the
retained $L_6$ from below; the enhanced half is unaffected.
The decrease of the paid loss relative to the value
in \eqref{num:lower-sum} is exactly
\[
 0.002372094-\lambda_{\rm loss}=\frac{37209}{10^9}.
\]
Combining this with \eqref{num:gross-bound} gives the
raw increment
\[
 \Delta=\frac{558583+37209}{10^9}
       =\frac{37237}{62500000}.
\]
Allocate $\Delta/8$ to each of the finite-approximation
and coefficient-realization reserves.  The net retained
increment is therefore $3\Delta/4=0.000446844$.
All estimates refer to this same fixed set of rectangles,
and each density contribution appears once.  They give
$R_{\rm low}=R_{\rm base}+\Delta$ and
$B_{\rm low}=B_{\rm base}+\Delta/8$, as in
\eqref{eq:current-R} and \eqref{eq:current-B}.
